% =====================================================================
%  interfaces.tex
%  Interfaces de usuario (v3): arquitectura modular con branding oficial
%  RED PEPPER (#FF3131, #F64301, #1F6636), separación cliente/equipo,
%  fotografía WebP y componentes reutilizables accesibles WCAG 2.2 AA.
% =====================================================================

\FloatBarrier
\section{Interfaces de usuario}
\label{sec:interfaces}

\noindent{\small\textbf{Versión 4} \textperiodcentered{} plataforma web modular implementada en \texttt{prototipo/} (\texttt{core/}, \texttt{components/} y \texttt{apps/} sin dependencias externas) con identidad oficial \textbf{RED PEPPER — Sánguches, Salchipapas \& Más}, sujeta a validación final con el interesado.}\par\vspace{0.5em}

En esta versión se reorganiza el prototipo aplicando el sistema de diseño \textbf{Stitch} (\texttt{.stitch/DESIGN.md}) en \textbf{tres subgrupos operativos} diferenciados por dispositivo (\textbf{Celular}, \textbf{Tablet} y \textbf{PC}) y provistos con doble apariencia \textbf{Versión Blanca} (\texttt{\#FAF8F5}) y \textbf{Versión Oscura} (\texttt{\#11100E}): (1)~\textbf{Subgrupo 1 — Experiencia Cliente} (modelo \emph{PedidosYa}: lista compacta y barra inferior de carrito en celular, catálogo 3 columnas con portada editorial en PC); (2)~\textbf{Subgrupo 2 — Operaciones en Tienda} (mostrador \textbf{POS táctil $\ge 48$\,px} con calculadora de vuelto y selector móvil \emph{Carta\,|\,Comanda}, más tablero \textbf{Cocina KDS} de 3 columnas con línea de urgencia y bloqueo rápido de stock); y (3)~\textbf{Subgrupo 3 — Administración y Marca} (\emph{backoffice} PC con riel lateral, prorrateo P-14, NIIF 18, visor modal de comprobantes SUNAT y estudio interno de marca). Todos los tokens cumplen contraste WCAG 2.2 AA en ambos temas.

\paragraph{Evolución respecto a las versiones anteriores.} Respecto a v1--v3,
la v4 incorpora: (i)~estructuración explícita en 3 subgrupos con flujos reales por tipo de dispositivo; (ii)~conmutación instantánea entre Versión Blanca y Versión Oscura (\texttt{?theme=light|dark}); (iii)~herramientas operativas reales (calculadora de vuelto rápido POS, modal de comprobante SUNAT y barras proporcionales por canal); y (iv)~suite automatizada \texttt{scripts/qa\_verify.mjs} aprobando 70 combinaciones vista $\times$ tema $\times$ resolución (320\,px a 1920\,px), WCAG 2.2 AA y 4 pruebas E2E, junto con capturas separadas en \texttt{interfaces/celular/}, \texttt{interfaces/tablet/} y \texttt{interfaces/pc/}. Los importes iniciales corresponden a un escenario de demostración local y no sustituyen registros reales.

\subsection{Lado del cliente y atención en mostrador}
\label{subsec:interfaces-cliente-pos}

\begin{figure}[H]
  \centering
  \includegraphics[width=\linewidth,height=0.72\textheight,keepaspectratio]{01-cliente-carta-pedido.png}
  \caption{Interfaz 1 (v3): portal público del cliente — marca oficial
  RED PEPPER, carta digital con fotografía WebP, buscador, selección de
  canal, personalización de cremas y seguimiento en vivo (\emph{Pendiente}
  $\rightarrow$ \emph{En preparación} $\rightarrow$ \emph{Listo}, conforme a
  R-OBS-03 y R-ENT2-17).}
  \label{fig:cliente-carta}
\end{figure}

\begin{figure}[H]
  \centering
  \includegraphics[width=\linewidth,height=0.72\textheight,keepaspectratio]{02-pos-mostrador.png}
  \caption{Interfaz 2 (v3): terminal interno de mostrador POS táctil.
  Controles con altura mínima de 48\,px (RNF-05 / R-ENT2-02), edición previa
  a la confirmación (R-OBS-04), bloqueo de productos agotados (R-ENT2-05) y
  confirmación modal de cierre de jornada (RF-03).}
  \label{fig:pos}
\end{figure}

\subsection{Lado de cocina (KDS)}
\label{subsec:interfaces-cocina}

\begin{figure}[H]
  \centering
  \includegraphics[width=\linewidth,height=0.72\textheight,keepaspectratio]{03-cocina-comandas.png}
  \caption{Interfaz 3 (v3): monitor de comandas en cocina (KDS) y control
  rápido de stock. Refleja preparación al instante sin precocidos
  (E-OBS-02 / E-ENT1-03), transiciones autorizadas de estado (R-OBS-03) y
  cambio inmediato de disponibilidad sincronizado con Cliente y POS
  (R-ENT1-02).}
  \label{fig:cocina-kds}
\end{figure}

\subsection{Lado de administración y contabilidad}
\label{subsec:interfaces-admin}

\begin{figure}[H]
  \centering
  \includegraphics[width=\linewidth,height=0.75\textheight,keepaspectratio]{04-admin-dashboard-margen.png}
  \caption{Interfaz 4 (v3): panel gerencial de margen bruto diario y
  prorrateo (ACT-0001). Calcula el margen diario (RF-08 / R-PRE-04),
  permite alternar jornadas hábiles y feriados conforme a P-14
  (R-ENT2-07), desglosa ventas por canal (R-ENT2-17) y muestra el ranking
  de carta incluyendo productos con cero ventas (R-ENT1-04).}
  \label{fig:dashboard}
\end{figure}

\begin{figure}[H]
  \centering
  \includegraphics[width=\linewidth,height=0.75\textheight,keepaspectratio]{05-admin-abastecimiento-carta.png}
  \caption{Interfaz 5 (v3): gestión de abastecimiento, carta y auditoría
  inmutable. Permite registrar compras con validación (R-ENT2-16), corregir
  importes mediante modal accesible con motivo obligatorio (R-ENT2-08), dar
  de alta o editar precios preservando historial (R-PRE-12 / R-PRE-13) y
  consultar la bitácora (RF-20 / RNF-02).}
  \label{fig:admin-operaciones}
\end{figure}

\begin{figure}[H]
  \centering
  \includegraphics[width=\linewidth,height=0.76\textheight,keepaspectratio]{06-admin-cierre-niif18-sunat.png}
  \caption{Interfaz 6 (v3): cierre mensual NIIF 18, reglas tributarias y
  repositorio SUNAT. Separa Resultado Operativo, Gastos Financieros y BAI
  (RF-10 a RF-14), implementa la autorización de Jerry sobre reglas
  propuestas por el contador (R-ENT2-10), aplica bloqueo de periodo cerrado
  con modal de nota de ajuste (RNF-07 / R-ENT2-14) y filtra comprobantes
  con hash SHA-256 (RF-16 / RNF-08).}
  \label{fig:admin-contabilidad}
\end{figure}
